\begin{abstract}
When an LLM inference server slows down under load, it is often unclear where the time
goes. We measured vLLM serving a 0.5B-parameter model on one GPU in a series of
pre-registered experiments totaling 18{,}900 measured requests. Most runs were
observed from the client, the server's phase metrics, and GPU and CPU telemetry, and
the last experiments added a per-step scheduler log. Every prediction was published
before its data and scored after.

First, from 8 to 32 concurrent requests, past throughput saturation, the added latency
moves into each engine step: the share of latency before the first token falls
from 0.44 to 0.24 while median inter-token latency rises from 4.9 to 55\,ms. Second,
the scheduler's token budget decides where that delay appears. At 32 concurrent
requests, a larger budget gives a 7\,ms median and 18\% more throughput but a
169\,ms 99th percentile, and with the smallest budget queueing takes two thirds of
each request's latency. Third, with one request at a time, a fixed-size prompt chunk
costs up to 1.6$\times$ more at deeper context, in line with a count of attention
FLOPs, and a cost model fit on those data predicts the 99th-percentile inter-token
latency at 8 concurrent requests within 5\% (post hoc). The pre-registered test that a
step's cost tracks its scheduled tokens failed at 8 concurrent requests and passed at
32 with a timing measure we later found lags by one step.

Of 59 scored predictions, 40 held. We report the failures in full, including that
measurement flaw. All GPU time for the study cost \$2.59.
\end{abstract}
% src: counts from the ledger (reports/*); 18,900 = sum of each report's measured requests; numbers from Phase C, Phase D (budget-0512 C32: queue 1222.3 of mean E2E 1843.9 ms), and Milestone 5 reports
